\documentclass[12pt]{article}
\usepackage{amsmath, amssymb}

\begin{document}

\section*{Governing Equations for Applied Exercises}

% ---------------------------------------------------------
% Equations for Exercise 1: The Baseline J2 Effect
% ---------------------------------------------------------
\subsection*{Exercise 1: The Baseline $J_2$ Effect}
To calculate the fundamental $J_2$ perturbation rates, we first need the orbit's semi-latus rectum ($p$) and the mean motion ($n$). For circular orbits, $e = 0$, so $p = a$.

\begin{equation}
    p = a(1 - e^2)
\end{equation}

\begin{equation}
    n = \sqrt{\frac{\mu}{a^3}}
\end{equation}

The secular rates of change for the longitude of the ascending node ($\dot{\Omega}$) and the argument of periapsis ($\dot{\omega}$) are given by:

\begin{equation}
    \dot{\Omega} = -\frac{3}{2} n J_2 \left( \frac{R_{\oplus}}{p} \right)^2 \cos i
\end{equation}

\begin{equation}
    \dot{\omega} = \frac{3}{2} n J_2 \left( \frac{R_{\oplus}}{p} \right)^2 \left( 2 - \frac{5}{2} \sin^2 i \right)
\end{equation}


% ---------------------------------------------------------
% Equations for Exercise 2: Sun-Synchronous Orbit (SSO)
% ---------------------------------------------------------
\subsection*{Exercise 2: Designing a Sun-Synchronous Orbit (SSO)}
For an orbit to be Sun-Synchronous, the nodal precession rate ($\dot{\Omega}$) must exactly match the Earth's mean angular velocity around the Sun ($\dot{\alpha}_{\odot}$).

\begin{equation}
    \dot{\Omega}_{req} = \dot{\alpha}_{\odot} \approx 0.9856^{\circ}/\text{day} = 1.991 \times 10^{-7} \text{ rad/s}
\end{equation}

By rearranging the nodal precession equation, we can solve directly for the required inclination ($i$) given a specific semi-major axis:

\begin{equation}
    \cos i = \frac{\dot{\Omega}_{req}}{-\frac{3}{2} n J_2 \left( \frac{R_{\oplus}}{p} \right)^2}
\end{equation}


% ---------------------------------------------------------
% Equations for Exercise 3: Critical Inclination
% ---------------------------------------------------------
\subsection*{Exercise 3: Critical Inclination and Frozen Orbits}
To prevent the line of apsides from rotating (a "frozen orbit"), the apsidal drift rate must be zero ($\dot{\omega} = 0$). 

\begin{equation}
    \frac{3}{2} n J_2 \left( \frac{R_{\oplus}}{p} \right)^2 \left( 2 - \frac{5}{2} \sin^2 i \right) = 0
\end{equation}

Since the leading terms are non-zero, the geometric condition for critical inclination reduces to:

\begin{equation}
    2 - \frac{5}{2} \sin^2 i = 0 \quad \implies \quad \sin^2 i = \frac{4}{5}
\end{equation}


% ---------------------------------------------------------
% Equations for Exercise 4: Jupiter SSO
% ---------------------------------------------------------
\subsection*{Exercise 4: Interplanetary Perturbations - Jupiter SSO}
The physics remain identical to Earth, but the planetary constants ($R_J$, $\mu_J$, $J_{2,J}$) and the target precession rate change. The required nodal precession matches the Jovian year ($T_J$).

\begin{equation}
    \dot{\Omega}_{req, J} = \frac{2\pi}{T_{J}} \text{ rad/s}
\end{equation}

\begin{equation}
    \cos i_J = \frac{\dot{\Omega}_{req, J}}{-\frac{3}{2} n_J J_{2,J} \left( \frac{R_J}{p_J} \right)^2}
\end{equation}


% ---------------------------------------------------------
% Equations for Exercise 5: Orbital Decay
% ---------------------------------------------------------
\subsection*{Exercise 5: Orbital Decay via Atmospheric Drag}
Atmospheric drag is a non-conservative force. Its effect on the spacecraft depends on the ballistic coefficient ($B$), which quantifies how susceptible the satellite is to drag.

\begin{equation}
    B = \frac{C_D A}{m}
\end{equation}

The time rate of change of the semi-major axis ($da/dt$) due to atmospheric drag in a nearly circular orbit is modeled as:

\begin{equation}
    \frac{da}{dt} = -n a^2 \rho B
\end{equation}
\textit{Where $\rho$ is the atmospheric density at the given altitude.}

\end{document}